\textbf{Phần 1. Trắc nghiệm lựa chọn}
\setcounterpageref{ex}{0}

\begin{ex} %[5D2N13]
Khi làm tròn số thập phân đến số tự nhiên gần nhất, ta so sánh chữ số ở hàng nào với số $5$?
\choice
{Hàng đơn vị}
{Hàng phần trăm}
{\True Hàng phần mười}
{Hàng phần nghìn}
\loigiai{
Khi làm tròn số thập phân đến số tự nhiên gần nhất, ta so sánh chữ số ở hàng phần mười với $5$.\\
\textbf{Đáp án của câu hỏi là C}
}
\end{ex}

\begin{ex} %[5D2N13]
Làm tròn số thập phân $31{,}2$ đến số tự nhiên gần nhất ta được
\choice
{$30$}
{\True $31$}
{$32$}
{$31{,}5$}
\loigiai{
Chữ số ở hàng phần mười là $2 < 5$ nên ta làm tròn xuống thành $31$.\\
\textbf{Đáp án của câu hỏi là B}
}
\end{ex}

\begin{ex} %[5D2H13]
Làm tròn số thập phân $31{,}75$ đến số tự nhiên gần nhất ta được
\choice
{\True $32$}
{$31$}
{$31{,}8$}
{$30$}
\loigiai{
Chữ số ở hàng phần mười là $7 > 5$ nên ta làm tròn lên thành $32$.\\
\textbf{Đáp án của câu hỏi là A}
}
\end{ex}

\begin{ex} %[5D2H13]
Làm tròn số thập phân $513{,}59$ đến số tự nhiên gần nhất ta được
\choice
{$513$}
{$510$}
{\True $514$}
{$513{,}6$}
\loigiai{
Chữ số ở hàng phần mười là $5$ nên ta làm tròn lên thành $514$.\\
\textbf{Đáp án của câu hỏi là C}
}
\end{ex}

\begin{ex} %[5D2H13]
Làm tròn số thập phân $0{,}806$ đến số tự nhiên gần nhất ta được kết quả là
\choice
{$0$}
{$0{,}8$}
{$2$}
{\True $1$}
\loigiai{
Chữ số ở hàng phần mười là $8 > 5$ nên ta làm tròn lên thành $1$.\\
\textbf{Đáp án của câu hỏi là D}
}
\end{ex}

\begin{ex} %[5D2H13]
Làm tròn số thập phân $2{,}52$ đến hàng phần mười ta được kết quả là
\choice
{$2{,}6$}
{$3{,}0$}
{\True $2{,}5$}
{$2{,}50$}
\loigiai{
Chữ số hàng phần trăm là $2 < 5$ nên ta làm tròn xuống được $2{,}5$.\\
\textbf{Đáp án của câu hỏi là C}
}
\end{ex}

\begin{ex} %[5D2H13]
Làm tròn số thập phân $3{,}25$ đến hàng phần mười ta được
\choice
{$3{,}2$}
{$3{,}0$}
{$4{,}0$}
{\True $3{,}3$}
\loigiai{
Chữ số hàng phần trăm là $5$ nên ta làm tròn lên được $3{,}3$.\\
\textbf{Đáp án của câu hỏi là D}
}
\end{ex}

\begin{ex} %[5D2H13]
Làm tròn số thập phân $6{,}324$ đến hàng phần trăm ta được
\choice
{\True $6{,}32$}
{$6{,}3$}
{$6{,}33$}
{$6{,}4$}
\loigiai{
Chữ số hàng phần nghìn là $4 < 5$ nên ta làm tròn xuống được $6{,}32$.\\
\textbf{Đáp án của câu hỏi là A}
}
\end{ex}

\begin{ex} %[5D2H13]
Làm tròn số thập phân $6{,}327$ đến hàng phần trăm ta được
\choice
{$6{,}32$}
{$6{,}30$}
{$6{,}4$}
{\True $6{,}33$}
\loigiai{
Chữ số hàng phần nghìn là $7 > 5$ nên ta làm tròn lên được $6{,}33$.\\
\textbf{Đáp án của câu hỏi là D}
}
\end{ex}

\begin{ex} %[5D2H13]
Làm tròn số $\pi = 3{,}141592\dots$ đến hàng phần trăm ta được
\choice
{$3{,}1$}
{\True $3{,}14$}
{$3{,}15$}
{$3{,}2$}
\loigiai{
Chữ số hàng phần nghìn của $\pi$ là $1 < 5$ nên ta làm tròn xuống được $3{,}14$.\\
\textbf{Đáp án của câu hỏi là B}
}
\end{ex}

\begin{ex} %[5D2V13]
Rô-bốt đo đường chéo của màn hình ti-vi $55\text{ inch}$ được kết quả là $139{,}7\text{ cm}$. Nếu làm tròn kết quả đó đến số tự nhiên gần nhất thì đường chéo ti-vi dài khoảng
\choice
{\True $140\text{ cm}$}
{$139\text{ cm}$}
{$139{,}5\text{ cm}$}
{$145\text{ cm}$}
\loigiai{
Số $139{,}7$ có chữ số hàng phần mười là $7 > 5$ nên làm tròn lên thành $140\text{ cm}$.\\
\textbf{Đáp án của câu hỏi là A}
}
\end{ex}

\begin{ex} %[5D2V13]
Một quả dưa hấu cân nặng chính xác là $2{,}52\text{ kg}$. Cô bán hàng đã làm tròn số cân nặng đến hàng phần mười. Hỏi quả dưa hấu cân nặng khoảng bao nhiêu ki-lô-gam?
\choice
{$2\text{ kg}$}
{\True $2{,}5\text{ kg}$}
{$2{,}6\text{ kg}$}
{$3\text{ kg}$}
\loigiai{
Số $2{,}52$ có chữ số hàng phần trăm là $2 < 5$ nên làm tròn đến hàng phần mười ta được $2{,}5\text{ kg}$.\\
\textbf{Đáp án của câu hỏi là B}
}
\end{ex}

\textbf{Phần 2. Trắc nghiệm đúng sai}
\setcounterpageref{ex}{0}

\begin{ex} %[5D2H13]
Cho ba số thập phân: $A = 9{,}345$; $B = 21{,}663$; $C = 0{,}4571$. Xét tính đúng sai của các khẳng định sau:
\choiceTF
{\True Làm tròn số $A$ đến hàng phần mười ta được $9{,}3$.}
{Làm tròn số $A$ đến hàng phần trăm ta được $9{,}34$.}
{\True Làm tròn số $B$ đến hàng phần trăm ta được $21{,}66$.}
{\True Làm tròn số $C$ đến hàng phần mười ta được $0{,}5$.}
\loigiai{
\begin{itemize}
    \item Số $A = 9{,}345$ có hàng phần trăm là $4 < 5 \Rightarrow$ làm tròn đến hàng phần mười là $9{,}3$.
    \item Số $A = 9{,}345$ có hàng phần nghìn là $5 \Rightarrow$ làm tròn đến hàng phần trăm là $9{,}35 \neq 9{,}34$.
    \item Số $B = 21{,}663$ có hàng phần nghìn là $3 < 5 \Rightarrow$ làm tròn đến hàng phần trăm là $21{,}66$.
    \item Số $C = 0{,}4571$ có hàng phần trăm là $5 \Rightarrow$ làm tròn đến hàng phần mười là $0{,}5$.
\end{itemize}
\begin{itemize}
    \item a) \textbf{Đúng}.
    \item b) \textbf{Sai}: Kết quả là $9{,}35$.
    \item c) \textbf{Đúng}.
    \item d) \textbf{Đúng}.
\end{itemize}
}
\end{ex}

\begin{ex} %[5D2V13]
Theo Tổ chức Y tế Thế giới (WHO), chiều cao và cân nặng chuẩn của bé trai $10$ tuổi lần lượt là $137{,}8\text{ cm}$ và $31{,}2\text{ kg}$; của bé gái $10$ tuổi là $138{,}6\text{ cm}$ và $31{,}9\text{ kg}$. Xét tính đúng sai của các khẳng định sau khi làm tròn đến số tự nhiên gần nhất:
\choiceTF
{\True Chiều cao chuẩn của bé trai làm tròn thành $138\text{ cm}$.}
{Cân nặng chuẩn của bé trai làm tròn thành $32\text{ kg}$.}
{\True Chiều cao chuẩn của bé gái làm tròn thành $139\text{ cm}$.}
{\True Cân nặng chuẩn của bé gái làm tròn thành $32\text{ kg}$.}
\loigiai{
\begin{itemize}
    \item $137{,}8$ có hàng phần mười là $8 > 5 \Rightarrow$ làm tròn thành $138\text{ cm}$.
    \item $31{,}2$ có hàng phần mười là $2 < 5 \Rightarrow$ làm tròn thành $31\text{ kg} \neq 32\text{ kg}$.
    \item $138{,}6$ có hàng phần mười là $6 > 5 \Rightarrow$ làm tròn thành $139\text{ cm}$.
    \item $31{,}9$ có hàng phần mười là $9 > 5 \Rightarrow$ làm tròn thành $32\text{ kg}$.
\end{itemize}
\begin{itemize}
    \item a) \textbf{Đúng}.
    \item b) \textbf{Sai}: Làm tròn thành $31\text{ kg}$.
    \item c) \textbf{Đúng}.
    \item d) \textbf{Đúng}.
\end{itemize}
}
\end{ex}

\textbf{Phần 3. Trả lời ngắn}
\setcounterpageref{ex}{0}

\begin{ex} %[5D2H13]
Làm tròn số thập phân $42{,}305$ đến số tự nhiên gần nhất.
\shortans{42}
\loigiai{
\textbf{Đáp số:} $42$\\
Chữ số hàng phần mười là $3 < 5$ nên làm tròn thành $42$.
}
\end{ex}

\begin{ex} %[5D2H13]
Làm tròn số thập phân $9{,}345$ đến hàng phần trăm. (Nhập kết quả dưới dạng số thập phân).
\shortans{9,35}
\loigiai{
\textbf{Đáp số:} $9{,}35$\\
Chữ số hàng phần nghìn là $5$ nên làm tròn lên thành $9{,}35$.
}
\end{ex}

\begin{ex} %[5D2H13]
Làm tròn số thập phân $21{,}663$ đến hàng phần mười. (Nhập kết quả dưới dạng số thập phân).
\shortans{21,7}
\loigiai{
\textbf{Đáp số:} $21{,}7$\\
Chữ số hàng phần trăm là $6 > 5$ nên làm tròn lên thành $21{,}7$.
}
\end{ex}

\begin{ex} %[5D2H13]
Làm tròn số $\pi = 3{,}141592\dots$ đến hàng phần mười. (Nhập kết quả dưới dạng số thập phân).
\shortans{3,1}
\loigiai{
\textbf{Đáp số:} $3{,}1$\\
Chữ số hàng phần trăm là $4 < 5$ nên làm tròn xuống thành $3{,}1$.
}
\end{ex}

\begin{ex} %[5D2V13]
Có bao nhiêu số thập phân có một chữ số ở phần thập phân lớn hơn $6$ mà khi làm tròn đến số tự nhiên gần nhất ta được kết quả bằng $7$?
\shortans{4}
\loigiai{
\textbf{Đáp số:} $4$\\
Các số đó là: $6{,}5;\; 6{,}6;\; 6{,}7;\; 6{,}8;\; 6{,}9$. Vì lớn hơn $6$ và có $1$ chữ số thập phân nên có $5$ số: $6{,}5; 6{,}6; 6{,}7; 6{,}8; 6{,}9$. Khi làm tròn đến số tự nhiên gần nhất, cả $5$ số này đều bằng $7$.
}
\end{ex}

\begin{ex} %[5D2V13]
Một mảnh vườn hình chữ nhật có chiều dài $12{,}4\text{ m}$ và chiều rộng $5{,}8\text{ m}$. Tính chu vi của mảnh vườn đó sau khi đã làm tròn mỗi kích thước đến số tự nhiên gần nhất theo đơn vị mét.
\shortans{36}
\loigiai{
\textbf{Đáp số:} $36$\\
Làm tròn chiều dài $12{,}4\text{ m} \rightarrow 12\text{ m}$.\\
Làm tròn chiều rộng $5{,}8\text{ m} \rightarrow 6\text{ m}$.\\
Chu vi mảnh vườn là: $(12 + 6) \times 2 = 36\text{ (m)}$.
}
\end{ex}
